\section{Expose one complete restricted-baseline macroblock}
Implement \texttt{h263\_emit\_macroblock}. Pass an eight-byte-aligned mapped
80-byte request in \texttt{\$a0}. Its fields are native-endian doublewords;
motion and predictor components are signed, while other fields are unsigned.
\begin{center}
  \tagpdfsetup{table/header-rows={1}}
  \begin{tabular}{rll}
    \textbf{Offset} & \textbf{Field}          & \textbf{Meaning}                           \\
    0               & Six-block level pointer & 768 tight bytes, two-byte aligned          \\
    8               & Writer-context pointer  & Earlier eight-byte-aligned 24-byte context \\
    16              & Picture type            & 0 I; 1 P                                   \\
    24              & INTRA selector          & 0 INTER; 1 INTRA                           \\
    32, 40          & Selected motion $x,y$   & Half-sample components, $-32$ through 31   \\
    48, 56          & Predictor $x,y$         & Same signed units and range                \\
    64              & Allow skip              & 0 coded; 1 permits eligible skipped coding \\
    72              & Reserved                & Zero
  \end{tabular}
\end{center}
The six arrays are consecutive 128-byte natural-order level blocks. All ordinary
levels lie in $[-127,127]$. For INTRA, index zero of every block instead holds
DC code one through 255, excluding 128. I-picture requests require INTRA. INTRA
requests require all four vector fields zero. INTER requests may use any admitted
selected/predictor components; picture geometry and predictor derivation belong
to the caller, which must supply a legal prediction for the actual picture.

The writer context retains its earlier one-MiB capacity bound. Request, levels,
context and output are disjoint, mapped and non-wrapping. Preserve request and
level bytes, context pointer/capacity, saved general registers and incoming stack
pointer. Supply an aligned stack with private space; vector registers are
caller-saved. This function changes no reconstructed picture or motion-history
array. The picture controller will update those only after accepting the encoded
macroblock or picture.

Malformed metadata or a forbidden level anywhere in the six blocks returns status
one/count zero. Metadata rejection precedes all level reads; level validation
reads are permitted after admitted metadata. No caller-buffer access or state
change occurs on either rejection. If the complete chosen representation does not
fit, return status two/count zero, also before caller-buffer access. Child failure
returns status three/count zero with caller bytes and context unchanged.

On success, append exactly one supported macroblock, retain bits outside it, commit
the used count and return status zero with appended bit count in \texttt{\$v1}.
The payload has no arbitrary padding or macroblock stuffing. Quantizer-change types,
four-vector modes and optional annex syntax are outside this entry. The selected
picture quantizer remains in force. Allow-skip zero requires coded output. With
allow-skip one, a zero-motion, zero-level INTER request in a P-picture may use
COD one or an equivalent coded INTER representation; both remain valid.

Use \path{construction/h263-macroblock-syntax/main.asm}. Complete the entry, then carry
in the completed coefficient emitter, its required code data and bit writer.
The provided composition helper accepts the earlier working format and retains
the macroblock caller and entry before its dependency marker:
\begin{verbatim}
python3 construction/h263-macroblock-syntax/compose.py \
    --source my-macroblock.asm --coefficients my-coefficients.asm \
    --output joined-macroblock.asm
python3 construction/h263-macroblock-syntax/prepare.py \
    --source joined-macroblock.asm --picture 0 --intra 1 \
    --mask 0 --used 3 --output intra-dc-only.asm
\end{verbatim}
Both helpers refuse an existing output. Preparation changes supplied levels and
controls while retaining all instructions and numeric tables. The default I-INTRA
observation has six DC codes 255 and no AC events. Its MCBPC uses one bit, CBPY
four, and the DC fields forty-eight: 53 appended bits. At initial used count
three, it ends on a byte boundary and prints 53 before exiting zero.

At \texttt{syntax\_caller\_return}, before the first caller instruction, output,
context and return registers are complete; the named status/bit-count fields
still hold their previous values. Inspect the decoded type, flags, motion and
all six recovered blocks as well as the final count. A correct bit count alone
cannot distinguish a swapped chroma flag from a correct payload.
